% Reviewer-facing draft for AE-4 / R1-4 (JACIII Jc26-0002).
% Section number is intentionally written 5.5 to match the rebuttal letter;
% adjust if the surrounding paper structure has shifted.
\subsection{Computational Performance and Real-Time Feasibility}
\label{sec:runtime}

We instrumented every STCM module with \texttt{time.perf\_counter}
boundaries inside the semantic-map builder node and replayed three
indoor rosbags offline (\texttt{meeting}, \texttt{livinglab}, and a
tuned variant \texttt{livinglab\_tuned}) on git SHA
\texttt{6880a5f}. Per-module mean / p50 / p95 latencies and the
effective per-bag rate are reported in Table~\ref{tab:runtime}; raw
samples and the per-scenario event counters used to interpret them
are archived under
\texttt{results/bench/runtime\_6880a5f.json}.

\noindent\textbf{Hardware.}
All numbers are measured on a single workstation: Intel Xeon
w7-3465X (28 cores / 56 threads), NVIDIA RTX 6000 Ada Generation
(48~GiB VRAM), 125~GiB system RAM, Ubuntu 22.04.5 LTS
(kernel 5.15.0-143-generic), ROS~2 Humble, CUDA toolkit 12.4,
PyTorch 2.4.0+cu121, NVIDIA driver 580.82.07. We disclose the host
explicitly because the perception backbone (GroundingDINO-T plus
MobileSAM) is GPU-bound; we discuss the implication for embedded
deployment at the end of the section.

\noindent\textbf{Critical-vs-asynchronous split.}
We separate the pipeline into modules that must close the
navigation loop in real time (\emph{critical}) and modules that may
run on independent threads and lag the navigation cycle without
stalling it (\emph{async}). FAST-LIO2 odometry, the per-frame
detection filter, the LiDAR-image association, and both growing-
neural-gas updaters are critical: their latency directly bounds the
rate at which the robot's pose and the topological graph can react
to new sensor data. GroundingDINO, MobileSAM, the temporal
score-fusion step, and the end-to-end per-frame pipeline are
asynchronous: they feed the semantic graph but do not gate Nav2
planning or controller cycles, which subscribe to the place graph
and instance graph through ROS topics.

\noindent\textbf{Critical-path budget.}
Summing the p95 latencies of the critical-path STCM modules at the
operating point used for evaluation gives
$0.3 + 13.2 + 22.5 + 22.1 \approx 58~\mathrm{ms}$ per detection cycle
(detection filter, LiDAR-image association, instance-GNG, place-GNG;
Table~\ref{tab:runtime}, averaged over three scenarios). FAST-LIO2
publishes \texttt{/fastlio2/lio\_odom} at $7.7$~Hz on average across
the three replays, well within the open-loop latency our place-graph
update can absorb ($\sim 22$~ms p95). The combined critical loop
therefore admits closed-loop operation at \emph{at least} $\sim 17$~Hz
on this hardware, which exceeds the rate at which an indoor mobility
aid is required to update its local plan~\cite{nav2}. No critical-
path module exceeded $25$~ms p95 in any scenario.

\noindent\textbf{Asynchronous semantic update.}
The slowest single module is GroundingDINO at
$164$~ms mean / $162$~ms p95, followed by MobileSAM at $63$~ms / $70$~ms
and temporal fusion at $40$~ms / $87$~ms. The end-to-end per-frame
pipeline takes $1.7$--$1.9~\mathrm{s}$ mean, which is dominated by
the perception cascade plus per-detection graph maintenance. Because
this stage is asynchronous, the practical consequence is a
\emph{semantic-update lag} of roughly one frame -- on the order of a
single second -- between the moment an object enters the camera
view and the moment the corresponding node is committed to the
shared graph. In a stationary or slow-moving assistive-navigation
setting this lag is acceptable: the obstacle layer underlying Nav2
is still served by the LiDAR pipeline at $7.7$~Hz, and the
semantic graph only needs to be fresh enough to ground a
language command, not to avoid collisions. This decoupling is the
load-bearing architectural claim of the system; we make it
explicit so a reader can judge it.

\noindent\textbf{Honesty caveats.}
We surface three known measurement biases rather than launder them:
(i) the \emph{Per-frame pipeline} sample includes the cheap
zero-detection / target-label-empty early-return paths
(70 / 350 frames on \texttt{meeting}, 21 / 54 on \texttt{livinglab}),
which deflates p50 and p95 relative to the cost of a frame that
actually triggers the full perception cascade -- the relevant event
counters (\texttt{zero\_detection\_frames},
\texttt{target\_label\_empty\_frames}) are recorded alongside every
\texttt{result.json} so a reviewer can recompute conditional
statistics; (ii) the \emph{LiDAR-image assoc} sample similarly
includes TF-lookup fast-fail paths, which is why p50 and p95 differ
by an order of magnitude
($1.4$~ms vs.\ $13.2$~ms; the failure count is
$113 / 458$ associations on \texttt{meeting}); and (iii) the first
GroundingDINO call carries CUDA-kernel compilation overhead that
inflates the mean very slightly above the steady-state cost. Nav2
planner and controller latencies are reported architecturally rather
than instrumented end-to-end; we treat the asynchronous-decoupling
guarantee as a design claim, not a measurement.

\noindent\textbf{Embedded-deployment outlook.}
The numbers in Table~\ref{tab:runtime} were obtained on a desktop
workstation with a high-end discrete GPU, not on the on-board
compute typically available on a wheelchair-class assistive
platform. Critical-path STCM modules are pure-CPU and would scale
within a factor of two on a laptop-grade Xeon; the perception
cascade is the dominant unknown. As future work we plan a parity
sweep on a Jetson AGX Orin -- the modal embedded target for indoor
service robots -- and to publish the resulting Table~\ref{tab:runtime}
delta. We do \emph{not} claim onboard real-time feasibility from
the present measurements; we claim only that, at the algorithmic
level, the critical-path budget leaves substantial headroom for
weaker hardware and the asynchronous design contains the perception
cost outside the navigation loop.

% Table D is generated by scripts/render_table_d.py; do not hand-edit.
% Generated by scripts/render_table_d.py — do not hand-edit
% AE-4 / R1-4: computational performance and real-time feasibility
\begin{table}[t]
\centering
\caption{Per-module runtime of the STCM pipeline (AE-4 / R1-4). Mean / p50 / p95 latencies averaged across 3 offline rosbag replays (livinglab, livinglab_tuned, meeting) on git SHA \texttt{6880a5f}; \emph{Eff.\,Hz} = sample count / bag duration. \emph{Critical} modules must close the navigation loop in real time; \emph{async} modules run on independent threads and do not block Nav2 / FAST-LIO2 (semantic graph updates lag detections by one cycle). Hardware: Intel(R) Xeon(R) w7-3465X (28c/56t) + NVIDIA RTX 6000 Ada Generation, 49140 MiB, 125 GiB RAM, Ubuntu 22.04.5 LTS kernel 5.15.0-143-generic, ROS 2 Humble, CUDA toolkit 12.4, PyTorch 2.4.0+cu121, NVIDIA driver 580.82.07.}
\label{tab:runtime}
\begin{tabular}{l l r r r r}
\toprule
Module & Class & Mean (ms) & p50 (ms) & p95 (ms) & Eff.\,Hz \\
\midrule
FAST-LIO2 odom (publish) & critical & -- & -- & -- & 7.7 \\
GroundingDINO detect & async & 164 & 157 & 162 & 0.9 \\
MobileSAM segment & async & 62.9 & 62.5 & 70.3 & 0.7 \\
Detection filter & critical & 0.2 & 0.2 & 0.3 & 0.9 \\
LiDAR-image assoc$^{\dagger}$ & critical & 4.3 & 1.4 & 13.2 & 1.5 \\
Instance-GNG update & critical & 9.0 & 11.1 & 22.5 & 1.3 \\
Place-GNG update & critical & 17.3 & 21.5 & 22.1 & 0.9 \\
Temporal fusion & async & 39.4 & 30.7 & 87.4 & 0.7 \\
Per-frame pipeline$^{\dagger}$ & async & 1693 & 1969 & 2468 & 0.9 \\
\bottomrule
\end{tabular}
\\[2pt]
\footnotesize $^{\dagger}$Sample includes cheap-path early returns (zero-detection / target-label-empty frames for \emph{Per-frame pipeline}; TF-lookup fast-fail for \emph{LiDAR-image assoc}); reader should consult \texttt{events.zero\_detection\_frames}, \texttt{target\_label\_empty\_frames}, and \texttt{pose\_failures} in \texttt{results/bench/runtime\_<sha>.json} for the denominator.
\end{table}

